% Options for packages loaded elsewhere
\PassOptionsToPackage{unicode}{hyperref}
\PassOptionsToPackage{hyphens}{url}
\PassOptionsToPackage{dvipsnames,svgnames,x11names}{xcolor}
\documentclass[
  11pt,
  a4paper,
]{article}
\usepackage{xcolor}
\usepackage[margin=25mm]{geometry}
\usepackage{amsmath,amssymb}
\setcounter{secnumdepth}{-\maxdimen} % remove section numbering
\usepackage{iftex}
\ifPDFTeX
  \usepackage[T1]{fontenc}
  \usepackage[utf8]{inputenc}
  \usepackage{textcomp} % provide euro and other symbols
\else % if luatex or xetex
  \usepackage{unicode-math} % this also loads fontspec
  \defaultfontfeatures{Scale=MatchLowercase}
  \defaultfontfeatures[\rmfamily]{Ligatures=TeX,Scale=1}
\fi
\usepackage{lmodern}
\ifPDFTeX\else
  % xetex/luatex font selection
\fi
% Use upquote if available, for straight quotes in verbatim environments
\IfFileExists{upquote.sty}{\usepackage{upquote}}{}
\IfFileExists{microtype.sty}{% use microtype if available
  \usepackage[]{microtype}
  \UseMicrotypeSet[protrusion]{basicmath} % disable protrusion for tt fonts
}{}
\makeatletter
\@ifundefined{KOMAClassName}{% if non-KOMA class
  \IfFileExists{parskip.sty}{%
    \usepackage{parskip}
  }{% else
    \setlength{\parindent}{0pt}
    \setlength{\parskip}{6pt plus 2pt minus 1pt}}
}{% if KOMA class
  \KOMAoptions{parskip=half}}
\makeatother
\usepackage{longtable,booktabs,array}
\newcounter{none} % for unnumbered tables
\usepackage{calc} % for calculating minipage widths
% Correct order of tables after \paragraph or \subparagraph
\usepackage{etoolbox}
\makeatletter
\patchcmd\longtable{\par}{\if@noskipsec\mbox{}\fi\par}{}{}
\makeatother
% Allow footnotes in longtable head/foot
\IfFileExists{footnotehyper.sty}{\usepackage{footnotehyper}}{\usepackage{footnote}}
\makesavenoteenv{longtable}
\setlength{\emergencystretch}{3em} % prevent overfull lines
\providecommand{\tightlist}{%
  \setlength{\itemsep}{0pt}\setlength{\parskip}{0pt}}
\setlength{\parindent}{0pt}
\setlength{\parskip}{\baselineskip}
\usepackage{bookmark}
\IfFileExists{xurl.sty}{\usepackage{xurl}}{} % add URL line breaks if available
\urlstyle{same}
\hypersetup{
  colorlinks=true,
  linkcolor={blue},
  filecolor={Maroon},
  citecolor={Blue},
  urlcolor={blue},
  pdfcreator={LaTeX via pandoc}}

\author{}
\date{}

\begin{document}

\section{Response to the Associate Editor and
Reviewers}\label{response-to-the-associate-editor-and-reviewers}

\textbf{Manuscript:} \emph{FlowMOP: An Automated Flow Cytometry Time,
Debris, and Doublet Removal Tool}

The Associate Editor and reviewers raised several overlapping concerns.
To avoid repeating the same response while ensuring that every point is
addressed, closely related comments are reproduced verbatim and answered
together below. A complete comment-coverage index is provided at the end
of this response.

We thank the Associate Editor and reviewers for their considered
feedback and believe that the requested amendments have significantly
improved the manuscript. We hope that the revised manuscript now
addresses the Associate Editor\textquotesingle s and
reviewers\textquotesingle{} concerns.

\subsection{Data Leakage: Correction to the Comparative
Benchmark}\label{data-leakage-correction-to-the-comparative-benchmark}

During revision, we identified a data-leakage error in the original
competitor benchmark: the source-label field retained for scoring had
inadvertently been available to PeacoQC and FlowCut. The synthetic FCS
files deliberately retained \texttt{SampleIDInt} so that event-level
source truth remained attached for scoring, but the original PeacoQC and
FlowCut benchmark paths did not remove this numeric field from the
channels available for quality assessment. FlowMOP did not have this
issue because its fluorescence-channel selection explicitly excluded
channel names containing \texttt{sample}, together with Time and scatter
channels.

We reran the full benchmark with the source label excluded from all
quality-control inputs and used it only for scoring. The correction
reinforced rather than reversed the comparative findings. FlowMOP had
higher sensitivity than FlowCut for Segment at both tested bin sizes and
for the 5000-event Bimix and Trimix benchmarks. At 5000 events, FlowMOP
also had higher specificity than both competitors across Segment, Bimix,
and Trimix; at 2000 events, it retained higher specificity than PeacoQC
in Bimix and Trimix without a significant sensitivity disadvantage and
had higher sensitivity and specificity than FlowCut in Segment. Figure
2, its violin distributions, p-values, significance brackets, and all
reported comparisons now use these corrected outputs.

This correction is independent of the specific reviewer comments
addressed below and is reported for transparency.

\subsection{Associate Editor\textquotesingle s General
Assessment}\label{associate-editors-general-assessment}

\begin{quote}
Novel tools to expedite flow cytometry data analysis are always welcome,
and I commend the authors for their effort in introducing these new
methods. However, as currently submitted, the manuscript requires major
revision. The reviewers have identified numerous areas requiring
improvement, and I encourage the authors to address their comments
carefully and comprehensively. In addition to the
reviewers\textquotesingle{} critiques, I offer several observations of
my own that I believe will strengthen the manuscript\textquotesingle s
suitability for publication.
\end{quote}

\textbf{Response:} The revisions are organized around the algorithmic
motivation, comparative mechanism, computational benchmarking,
interpretation of expert evaluations, methodological clarity, and the
limits of the current approach. Detailed responses follow.

\subsection{Combined Comment 1 --- Algorithmic Motivation, Novelty,
Related Tools, and Validation
Framing}\label{combined-comment-1--algorithmic-motivation-novelty-related-tools-and-validation-framing}

\textbf{Raised by:} Associate Editor; Reviewer 2, Comments 1, 6, and 8.

\subsubsection{Associate Editor --- Algorithmic Motivation and
Theoretical
Justification}\label{associate-editor--algorithmic-motivation-and-theoretical-justification}

\begin{quote}
The authors acknowledge the existence of other algorithms that perform
substantially similar analyses to the components of FlowMop, yet they do
not adequately motivate what deficiencies in the existing algorithms
their new approach is designed to address. A compelling framing would
follow a structure such as: we demonstrated that PeacoQC and FlowCut
fail under realistic conditions X, Y, and Z; upon analysis, we
determined this is attributable to specific aspects of how those
algorithms operate; accordingly, we designed our algorithm to overcome
these limitations. No such explanation is currently provided, nor is
there a clear justification for why the particular algorithmic choices
were made.
\end{quote}

\subsubsection{Reviewer 2, Comment 1}\label{reviewer-2-comment-1}

\begin{quote}
In the abstract, the authors claim ``FlowMOP is the first automated
approach capable of debris cleaning''. However, there are methods like
GateNet, 10.1002/cyto.a.20531, and UNITO that can do automatic gating
and probably can be used to gate out debris. There is even a review for
this purpose (10.3389/fimmu.2015.00380).
\end{quote}

\subsubsection{Reviewer 2, Comment 6}\label{reviewer-2-comment-6}

\begin{quote}
The introduction did not introduce available related tools.
\end{quote}

\subsubsection{Reviewer 2, Comment 8}\label{reviewer-2-comment-8}

\begin{quote}
The authors claimed, ``Conversely, this paper seeks to compare
performance not only through traditional human-expert defined standards,
but also the development of bespoke data generated explicitly for
pre-processing validation.'' In fact, both simulation data and read data
are commonly used for new method validation.
\end{quote}

\subsubsection{Response}\label{response}

We now frame FlowMOP as a training-free preprocessing workflow that
combines automated time, debris, and doublet removal. We revised the
Introduction to distinguish automated population-gating tools,
time-quality-control tools, and FlowMOP\textquotesingle s intended role
as an integrated preprocessing workflow. We removed the claim that
FlowMOP is the first automated approach capable of debris cleaning. The
revised Introduction discusses FlowCut and PeacoQC as time-dependent
quality-control approaches and GateNet and UNITO as broader automated
population-gating approaches.

We have also revised the description of the three FlowMOP modules. Time
gating is based on positive-population fluorescence summaries across
acquisition order, debris gating derives an FSC-A threshold from
cross-parameter positive-population structure, and doublet gating uses
FSC-A/FSC-H and SSC-A/SSC-H ratio structure. These algorithmic choices
are now distinguished from computational implementation choices.

In response to Reviewer 2\textquotesingle s observation that simulated
and real data are both commonly used in method validation, we now
describe the synthetic samples specifically as event-labelled technical
controls developed to provide objective ground truth for flow-cytometry
preprocessing tasks.

\subsubsection{Changes made}\label{changes-made}

\begin{itemize}
\item
  We removed the claim that FlowMOP is the first automated approach
  capable of debris cleaning and added the following clarification:

  \begin{quote}
  ``Several tools automate cytometry population identification,
  including the model-based flowClust method {[}13{]}, the
  neural-network method GateNet {[}14{]}, and the bivariate-segmentation
  framework UNITO {[}15{]}; automated gating approaches are reviewed
  elsewhere {[}16{]}. FlowMOP is a Python-based, training-free
  preprocessing workflow that combines automated time-gating, debris
  removal, and doublet exclusion in a single headless tool. FlowCut and
  PeacoQC provide the most direct comparison for its time-gating
  component.''
  \end{quote}
\item
  We reframed the synthetic-data contribution as follows:

  \begin{quote}
  ``The synthetic datasets provide event-level labels for estimating
  sensitivity and specificity in preprocessing tasks where real ground
  truth is otherwise difficult to define.''
  \end{quote}
\item
  We added a concise description of the three algorithmic components to
  the Figure 1 legend:

  \begin{quote}
  ``A) FlowMOP selects valid parameters using positive-peak detection,
  generates time-binned fluorescence summaries, applies two smoothing
  resolutions, and performs robust outlier rejection across acquisition
  order. B) FlowMOP applies the same valid-parameter check, derives a
  candidate FSC-A threshold from each eligible parameter's positive
  events, and uses the median candidate as the final FSC-A gate. C)
  FlowMOP identifies doublets from inflection points in FSC-A/FSC-H and
  SSC-A/SSC-H ratio histograms, with a fixed-ratio fallback.''
  \end{quote}
\end{itemize}

\textbf{Location:} Introduction; Results, ``Algorithmic design,'' ``Time
Gating,'' ``Debris Gating,'' and ``Doublet Gating''; Figure 1 legend.

\subsection{Combined Comment 2 --- Computational Benchmarking and
Within-Sample
Parallelisability}\label{combined-comment-2--computational-benchmarking-and-within-sample-parallelisability}

\textbf{Raised by:} Associate Editor; Reviewer 1, Comment 1; Reviewer 2,
Comment 7.

\subsubsection{Reviewer 1, Comment 1 --- Technical Infrastructure and
Performance
Claims}\label{reviewer-1-comment-1--technical-infrastructure-and-performance-claims}

\begin{quote}
The adoption of a Dask-based Python framework is a highly commendable
architectural choice. Flow cytometry analysis has historically
under-utilized distributed compute resources, and providing a modern,
scalable alternative to R-based packages is a significant contribution
to the ``Big Data'' era of spectral cytometry. However, the manuscript
currently lacks empirical data to support the claims of improved runtime
or memory efficiency, while other contemporary algorithms explicitly
highlight their scalability for large datasets.

Recommendation: To make a compelling case for FlowMOP's adoption, the
authors could include benchmarking figures (e.g., execution time and
peak memory usage) comparing FlowMOP against PeacoQC and FlowCut across
datasets of increasing file count and/or size (\textasciitilde10\^{}3 to
\textasciitilde10\^{}7 events or some {[}Event x Parameter{]} matrix
varying each). Another point of comparison may be in the distributed
compute scenario vs. more traditional local compute resources.

Example Performance Benchmarking Table:

{\def\LTcaptype{none} % do not increment counter
\begin{longtable}[]{@{}
  >{\raggedright\arraybackslash}p{(\linewidth - 8\tabcolsep) * \real{0.1920}}
  >{\raggedright\arraybackslash}p{(\linewidth - 8\tabcolsep) * \real{0.1920}}
  >{\raggedright\arraybackslash}p{(\linewidth - 8\tabcolsep) * \real{0.1920}}
  >{\raggedright\arraybackslash}p{(\linewidth - 8\tabcolsep) * \real{0.1920}}
  >{\raggedright\arraybackslash}p{(\linewidth - 8\tabcolsep) * \real{0.1920}}@{}}
\toprule\noalign{}
\begin{minipage}[b]{\linewidth}\raggedright
Dataset Size (Events)
\end{minipage} & \begin{minipage}[b]{\linewidth}\raggedright
Metric
\end{minipage} & \begin{minipage}[b]{\linewidth}\raggedright
FlowMOP (Python/Dask)
\end{minipage} & \begin{minipage}[b]{\linewidth}\raggedright
PeacoQC (R)
\end{minipage} & \begin{minipage}[b]{\linewidth}\raggedright
FlowCut (R)
\end{minipage} \\
\midrule\noalign{}
\endhead
\bottomrule\noalign{}
\endlastfoot
10\^{}4 & Execution Time (s) & & & \\
& Peak RAM (MB) & & & \\
10\^{}5 & Execution Time (s) & & & \\
& Peak RAM (MB) & & & \\
...etc. & & & & \\
\end{longtable}
}
\end{quote}

\subsubsection{Associate Editor --- Algorithmic vs. Implementation
Advantages}\label{associate-editor--algorithmic-vs-implementation-advantages}

\begin{quote}
Reviewer 1 raises several points regarding the use of the Dask
framework, and I would like to extend that line of inquiry.
Specifically, I would ask the authors to clarify whether there is
something inherent to the FlowMop algorithm that makes it particularly
well-suited to parallelization via Dask --- as distinct from the
implementation itself. This distinction matters considerably. Given
current development tools, a competent programmer with AI tooling could
likely port the existing FlowCut or PeacoQC code --- neither of which is
especially lengthy --- to a Dask-based parallel implementation in a
matter of hours. If that is the case, it raises the question of whether
a substantial portion of FlowMop\textquotesingle s reported advantage
derives from the implementation rather than from any intrinsic
algorithmic innovation, which would significantly diminish the novelty
of the contribution. I would encourage the authors to focus their
comparative analysis on algorithmic advantages, particularly in cases
where existing algorithms present genuine structural barriers to
parallelization.
\end{quote}

\subsubsection{Reviewer 2, Comment 7}\label{reviewer-2-comment-7}

\begin{quote}
The authors should introduce ``Dask'' before using the ``Dask-based
framework'' in the introduction. Please keep in mind, most readers
probably have no experience using Python at all. I use Python a lot
myself. Still, I am not familiar with this ``Dask'' thing. Also, the
Dask-based design is about calculation efficiency, instead of accuracy.
Readers will be more interested in which algorithm is used for automatic
gating, instead of how to speed up calculation. The authors did not
really introduce the algorithm for gating in the introduction. But they
use a whole paragraph on the Dask-based design.
\end{quote}

\subsubsection{Response}\label{response-1}

We evaluated computational performance using matched 36-channel FCS
inputs containing 10,000, 100,000, 300,000, 1,000,000, and 2,000,000
events. FlowMOP, PeacoQC, and FlowCut were run on the same clone-based
inputs. Optional plotting, reporting, and output generation were
disabled where supported, and FlowMOP\textquotesingle s debris and
doublet modules were disabled so that the shared time-gating task was
timed fairly. Each condition included one warm-up and three measured
repeats. Runtime and peak resident memory are reported as mean ± SD in
Table 1. At 2,000,000 events, FlowMOP was approximately 2.5-fold faster
than PeacoQC and 3.5-fold faster than FlowCut while using approximately
32\% and 45\% of their respective peak RAM.

The benchmark used local non-distributed execution. Testing active Dask
execution showed that its overhead did not improve this workload, so
Dask was inactive for the reported benchmark and removed from the
manuscript\textquotesingle s novelty framing.

Among the three evaluated decision structures, only FlowMOP is
inherently suited to independent within-sample channel parallelisation.
Once shared acquisition bins have been defined, each eligible channel
independently establishes its fluorescence reference, produces a
fixed-shape time-bin summary, and applies smoothing and MAD filtering.
Cross-channel coordination occurs only when the resulting bin flags are
combined in the final parameter vote. This creates coarse, regular
channel-level tasks with a single final reduction.

PeacoQC must instead reconcile peak configurations across bin-channel
combinations before its MAD, Isolation Tree, and consecutive-bin
decisions. FlowCut combines segment- and channel-level statistics
through adjacent-segment and file-wide comparisons, contiguous-region
decisions, and an optional flagged-file rerun. Both methods can process
separate files in parallel, and some internal suboperations could be
parallelised, but their complete within-sample decision pipelines
contain intermediate dependencies that prevent the same independent
channel-wise reduction without changing their decision rules.
Reimplementation in Python or another scheduling framework would not
remove those structural dependencies.

FlowMOP\textquotesingle s within-sample parallelisability is therefore
an inherent property of its decision structure rather than its
programming language. The reported runtime benchmark used local
non-distributed execution, so it measures the performance of the tested
implementations rather than an empirical distributed-computing speedup.

\subsubsection{Changes made}\label{changes-made-1}

\begin{itemize}
\item
  We removed Dask from the manuscript\textquotesingle s novelty and
  accuracy framing and specified the benchmark configuration as follows:

  \begin{quote}
  ``For fair timing of the shared time-gating task, FlowMOP was run
  using local non-distributed execution, with debris and doublet removal
  disabled and annotated output FCS writing disabled. PeacoQC and
  FlowCut were run with optional plotting, reporting, and output
  generation disabled where supported.''
  \end{quote}
\item
  We state FlowMOP\textquotesingle s contribution and identify the most
  direct time-gating comparators:

  \begin{quote}
  ``FlowMOP is a Python-based, training-free preprocessing workflow that
  combines automated time-gating, debris removal, and doublet exclusion
  in a single headless tool. FlowCut and PeacoQC provide the most direct
  comparison for its time-gating component.''
  \end{quote}
\item
  We retained a concise operational description in the Methods:

  \begin{quote}
  ``FlowMOP accepts \texttt{.csv}, \texttt{.fcs}, and Parquet files. It
  reduces event-level measurements into time-bin and channel summaries,
  applies smoothing and robust outlier detection, and combines the
  resulting flags through parameter voting.''
  \end{quote}
\item
  We placed the architectural comparison and its relationship to the
  computational benchmark in the Discussion:

  \begin{quote}
  ``These measurements were obtained using local non-distributed
  execution and therefore demonstrate the performance of the tested
  implementations, not a distributed-computing speedup.''
  \end{quote}

  \begin{quote}
  ``Among the three evaluated decision structures, only FlowMOP exposes
  an inherently channel-parallel within-sample computation. After shared
  acquisition bins are defined, each eligible fluorescence channel
  independently establishes its reference, produces a fixed-shape
  time-bin summary, and applies smoothing and MAD filtering;
  cross-channel coordination occurs only in the final parameter vote.
  PeacoQC instead performs data-dependent peak identification and
  reconciliation across bin-channel combinations, while FlowCut applies
  adjacent-segment, contiguous-region, file-wide, and conditional rerun
  decisions {[}5,9{]}. Their suboperations and separate files can be
  parallelised, but their complete within-sample decision pipelines
  cannot be expressed as the same independent channel-wise reduction
  without changing their decision rules. Porting either implementation
  to another language or scheduling framework would therefore not remove
  these structural dependencies. This structural distinction, together
  with FlowMOP's earlier data reduction, fewer full-data passes, and
  smaller intermediate state, is consistent with its lower runtime and
  memory use in our benchmarks, although the benchmark does not
  independently establish causation.''
  \end{quote}
\item
  We added matched runtime and peak-memory testing across five event
  counts:

  \begin{quote}
  ``Computational scalability was evaluated using clone-based real-FCS
  scaling. A representative FCS file was subsampled/replicated to
  matched event counts of 10,000, 100,000, 300,000, 1,000,000, and
  2,000,000 events while preserving the original 36-channel structure.
  FlowMOP, PeacoQC, and FlowCut were run on the same generated inputs
  for each size.''
  \end{quote}
\item
  We reported the principal computational result as follows:

  \begin{quote}
  ``The computational benchmark demonstrates that FlowMOP provides speed
  gains with lower peak RAM usage at larger event counts. This is most
  evident from 300,000 events onward, where FlowMOP had both the fastest
  mean runtime and lowest peak memory use. At 2,000,000 events, FlowMOP
  was approximately 2.5-fold faster than PeacoQC and 3.5-fold faster
  than FlowCut, while using approximately 32\% of PeacoQC's peak RAM and
  45\% of FlowCut's peak RAM.''
  \end{quote}
\end{itemize}

\textbf{Location:} Methods, ``Algorithmic design'' and ``Computational
scalability''; Results, ``Computational scalability''; Table 1;
Discussion, ``Synthetic Sample Time Gating.''

\subsection{Combined Comment 3 --- Downstream Biological Validation,
Dataset Breadth, and the Cost of
Errors}\label{combined-comment-3--downstream-biological-validation-dataset-breadth-and-the-cost-of-errors}

\textbf{Raised by:} Associate Editor; Reviewer 1, Comments 3 and 4;
Reviewer 2, general assessment and Comment 10.

\subsubsection{Associate Editor --- Downstream Biological
Impact}\label{associate-editor--downstream-biological-impact}

\begin{quote}
Alternatively --- or in addition --- the authors could evaluate the
downstream biological impact of each gating strategy. By substituting
FlowMop, FlowCut, and PeacoQC gates for the time and debris gating steps
while retaining the human gating strategy for all other steps, one could
directly assess whether differences in automated gating performance have
any meaningful effect on the biological conclusions drawn from the data.
Ultimately, the relevant question for end users is not whether an
algorithm removes a few percent more or fewer debris events, but whether
doing so materially affects downstream analysis.
\end{quote}

\subsubsection{Reviewer 1, Comment 3 --- Biological Dataset
Breadth}\label{reviewer-1-comment-3--biological-dataset-breadth}

\begin{quote}
Methodological Robustness: Utilizing a single expert per tissue type
limits the statistical power of the ranking methodology. It prevents the
reader from seeing the intra-tissue consensus---how consistently
multiple experts on the same tissue agree on the rank-order of methods
used on a specific sample. What it currently highlights is how a single
tissue expert has a unique perspective of the correct gating for a given
tissue and inter-operator gating variability; not algorithmic gating
superiority. Other improvements could be adding different processing
conditions (e.g. fixation) or more variation in tissue used (e.g.
include human PBMC or other species\textquotesingle{} tissues).
\end{quote}

\subsubsection{Reviewer 1, Comment 4 --- Discussion of Error ``Costs''
and Sensitivity
Trade-offs}\label{reviewer-1-comment-4--discussion-of-error-costs-and-sensitivity-trade-offs}

\begin{quote}
The manuscript notes that PeacoQC exhibits higher sensitivity but lower
specificity compared to FlowMOP. This trade-off is central to clinical
discovery and warrants a deeper discussion.

Discussion Point: The authors should explicitly address the biological
``cost'' of these errors. For instance, is it more detrimental to the
study\textquotesingle s integrity to leave in a small transient artifact
or to accidentally remove a rare, clinically significant biological
population?
\end{quote}

\subsubsection{Reviewer 2 --- Difficult Biological Samples and Tumour
Validation}\label{reviewer-2--difficult-biological-samples-and-tumour-validation}

\begin{quote}
The authors developed An Automated Flow Cytometry Time, Debris, and
Doublet Removal Tool. This can be a useful tool for cytometry. However,
the dataset used in the paper are relatively ``easy'' cases. I recommend
the authors to challenge the FlowMOP with more difficult but common
samples like digested tumor sample. Performance against human expert is
needed.
\end{quote}

\subsubsection{Reviewer 2, Comment 10}\label{reviewer-2-comment-10}

\begin{quote}
The study did not validate the FlowMOP on tumor datasets, which is a
common sample type for flow cytometry but difficult to perform debris
removal. I highly suggest the authors to include a tumor dataset, even
if the performance is not optimal, the readers can know whether the
method can be applied in their sample or not.
\end{quote}

\subsubsection{Response}\label{response-2}

We elected to conduct further biological validation, and therefore added
an experiment with 36 matched human PBMC samples and three human
tumour-liver samples. These analyses quantify both population recovery
and biological composition.

We agree with Reviewer 1 that the original single-expert-per-tissue
design cannot establish within-tissue expert consensus or support claims
of algorithmic superiority. We therefore present it as an
expert-preference evaluation, retain the complete analysis in the
Supplementary Information, and report its principal comparative findings
in the main Results, as detailed in Combined Comment 6. To broaden the
biological validation, we added the PBMC and tumour datasets described
here; we did not add a separate fixation or cross-species experiment.

For the PBMC analysis, we compared the time-gating methods while holding
the downstream expert-defined singlet, debris, Live CD45+, and lineage
gates fixed. We also evaluated the debris and doublet modules separately
against matched expert inputs in which only the relevant cleaning step
differed, and compared the complete Time + Debris + Doublet workflow
with the expert-combined workflow.

B cells were defined as Live CD45+ CD3−CD19+ events, T cells as Live
CD45+ CD3+CD19− events, and NKT cells as Live CD45+ CD19−CD3+CD56+
events. We report both counts and frequencies normalized to their
corresponding matched Raw values (Raw = 1). Before normalization, Live
CD45+ frequency was calculated relative to Live cells and B-, T-, and
NKT-cell frequencies relative to Live CD45+ cells. One sample was
excluded from NKT analyses because no NKT cells were detected (n = 35);
all other analyses used n = 36. Two-sided paired tests were
Holm-adjusted separately within endpoint and outcome metric.

The complete numerical results and statistical comparisons are reported
in the revised Results and Supplementary Data. In summary, the
time-gating methods retained different numbers of cells: FlowMOP and
FlowCut retained more than Expert Manual, whereas PeacoQC retained
fewer, but none of the evaluated downstream population frequencies
differed among methods (Figure 5B--C). In the separate debris
comparison, FlowMOP retained fewer cells and produced a slightly lower
Live CD45+ frequency than Expert Manual, while B-, T-, and NKT-cell
frequencies did not differ. In the doublet comparison, FlowMOP again
retained slightly fewer cells and produced a slightly lower Live CD45+
frequency and a slightly higher T-cell frequency, while B- and NKT-cell
frequencies did not differ. After the complete Time + Debris + Doublet
workflow, FlowMOP and Expert Manual did not differ in aggregate
population counts or in B-, T-, or NKT-cell frequencies; FlowMOP
produced a modestly lower Live CD45+ frequency (Figure 6D--E). Manual
and FlowMOP preprocessing did not differ for any of the three evaluated
tumour endpoints (Figure 7).

The biological data alone cannot establish whether the lower retention
by Expert Manual and PeacoQC reflects more sensitive artifact removal or
whether the higher retention by FlowMOP and FlowCut reflects more
specific preservation of valid events. The Discussion therefore
interprets these findings alongside the source-labelled synthetic
results. FlowMOP\textquotesingle s strong combined
sensitivity-specificity profile and PeacoQC\textquotesingle s
comparatively lower specificity support the latter interpretation for
FlowMOP and suggest a tendency towards over-removal by PeacoQC. The
higher retention by FlowMOP is also consistent with the comparatively
blunt, block-based nature of manual time gating, which can remove valid
events together with an affected interval.

The trade-off between sensitivity and specificity reflects competing
biological risks rather than a purely statistical optimization. A more
permissive gate may protect genuine or rare populations while allowing
artifacts to remain; a more stringent gate may remove artifacts more
completely while also deleting valid biological events. No balance is
universally correct because the consequences of each error depend on the
intended downstream analysis. We therefore report both sensitivity and
specificity and interpret them alongside their effects on population
recovery and composition. Importantly, FlowMOP had the strongest overall
combined sensitivity-specificity profile in the source-labelled
synthetic time-gating benchmarks: its gains were not achieved simply by
exchanging one error type for the other.

These biological analyses better illustrate FlowMOP\textquotesingle s
capabilities and reinforce the findings demonstrated in the synthetic
datasets. Figure S8 additionally contrasts sequential manual gating with
FlowMOP\textquotesingle s parallel exclusion logic.

\subsubsection{Changes made}\label{changes-made-2}

\begin{itemize}
\item
  We added a 36-sample PBMC biological validation with matched-Raw
  counts and matched-Raw frequencies for Live CD45+, B, T, and NKT
  cells; frequencies were first calculated within their specified
  biological parent. Figure 5 reports time-cleaning representatives (A),
  count statistics (B), and frequency statistics (C); Figure 6 reports
  debris, doublet, and combined-cleaning representatives (A--C), count
  statistics (D), and frequency statistics (E).
\item
  We added a tumour validation comprising three tumour samples, with
  paired Raw, Manual, and FlowMOP comparisons of Live CD45+ cell count,
  B-cell frequency, and T-cell frequency. Figure 7 displays the matched
  gating plots and downstream comparisons.
\item
  We added a representative schematic showing sequential manual gates
  versus FlowMOP\textquotesingle s parallel exclusions and
  union/intersection logic (Figure S8).
\item
  We acknowledged that the single-expert-per-tissue rankings cannot
  estimate within-tissue consensus, retained the complete ranking
  analysis in the Supplementary Information, and added its principal
  comparative findings to the main Results. We broadened the validation
  with human PBMC and tumour datasets but did not add a separate
  fixation or cross-species experiment.
\item
  We added a direct discussion of competing error costs, including the
  debris-poor PBMC example:

  \begin{quote}
  ``The trade-off between sensitivity and specificity reflects competing
  biological risks rather than a purely statistical optimization. A more
  permissive gate may protect genuine or rare populations while allowing
  artifacts to remain; a more stringent gate may remove artifacts more
  completely while also deleting valid biological events. No balance is
  universally correct because the consequences of each error depend on
  the intended downstream analysis. In this context,
  FlowMOP\textquotesingle s synthetic time-gating results are notable
  because its performance gains were not achieved simply by exchanging
  one error type for the other. FlowMOP had the strongest overall
  combined sensitivity-specificity profile across the tested
  conditions.''
  \end{quote}
\end{itemize}

\textbf{Location:} Methods, ``Biological-validation analysis'' and
``Tumour biological-validation analysis''; Results, ``Biological
validation'' and ``Tumour biological validation''; Discussion,
``Biological validation'' and ``Other remarks''; Figures 5--7; Figure
S8.

\subsection{Combined Comment 4 --- Gating Mechanism, Smoothing, and
Parameter
Fairness}\label{combined-comment-4--gating-mechanism-smoothing-and-parameter-fairness}

\textbf{Raised by:} Associate Editor; Reviewer 1, Comment 5; Reviewer 2,
Comment 20.

\subsubsection{Associate Editor --- Algorithmic Motivation and
Theoretical
Justification}\label{associate-editor--algorithmic-motivation-and-theoretical-justification-1}

\begin{quote}
More troublingly, the authors themselves appear uncertain about why
their algorithm performs better in certain cases. In the Discussion,
they write that a performance difference ``may be attributed to
FlowMop\textquotesingle s \textquotesingle smoothing\textquotesingle{}
implementation.'' This speculation is unsatisfying --- can the authors
not test this directly by modifying the smoothing implementation and
evaluating whether the results change accordingly? While presenting
cherry-picked examples in which one algorithm outperforms another has
illustrative value, a rigorous justification of the theoretical
underpinnings of the algorithmic design is essential. Without it,
readers cannot reasonably assess which sample types and experimental
conditions are expected to favor FlowMop over existing alternatives.
\end{quote}

\subsubsection{Reviewer 1, Comment 5 --- Parametric Sensitivity
Analysis}\label{reviewer-1-comment-5--parametric-sensitivity-analysis}

\begin{quote}
To strengthen the claim of ``superiority,'' it is essential to
demonstrate that the results are not merely a function of suboptimal
default settings for the competing algorithms.

Recommendation: A supplemental figure showing a sensitivity
analysis---where the parameters of PeacoQC and FlowCut are tuned across
a range---would provide a more ``apples-to-apples'' comparison and
confirm that FlowMOP's performance gains are algorithmic rather than
parametric.
\end{quote}

\subsubsection{Reviewer 2, Comment 20}\label{reviewer-2-comment-20}

\begin{quote}
The authors say, ``Finally, FlowMOP also novelly employs the use of
smoothing at multiple resolutions to ensure multiple types of
abherrations are detected''. What are the ``multiple types of
abherrations''? please clarify.
\end{quote}

\subsubsection{Response}\label{response-3}

We now seek to address the mechanisms underlying
FlowMOP\textquotesingle s superior performance under the primary fixed
comparison settings in the tested Segment, Bimix, and Trimix scenarios.
Our hypotheses arose from the different signals used by the methods and
from the pattern observed in the original benchmark. FlowMOP uses
acquisition order to define time bins but bases its removal decision on
fluorescence summaries of globally defined positive populations; Time
itself is excluded from the quality variables. FlowCut, by contrast, can
respond to local acquisition-density structure. This distinction was
particularly relevant because the Segment files retained stronger
source-linked Time-density structure and showed the clearest FlowCut
performance loss, whereas Time was approximately normalized in the
constructed Bimix and Trimix files. We therefore hypothesised that
FlowCut\textquotesingle s relative loss arose, at least in part, from
sensitivity to local acquisition rate rather than to the source-linked
fluorescence defect itself.

To test this hypothesis, we conducted a matched Time-only mechanism
analysis. This involved changing only the Time channel while preserving
fluorescence, scatter, source labels, and event order across 30
source-labelled Bimix, Trimix, and Segment files. We found that FlowMOP
and PeacoQC were both unchanged under source-linked and random Time
warping, whereas FlowCut\textquotesingle s removal behaviour changed
when local Time density was altered. The clearest loss occurred in
Segment inputs under source-linked Time warping.

PeacoQC\textquotesingle s documentation identifies acquisition-bin size
as a trade-off between the accuracy of within-bin density estimation,
the number of bins available for evaluating signal stability, and the
number of events affected when a bin is removed {[}5{]}. The Bimix and
Trimix benchmarks contain short, interspersed source-defined intervals.
FlowMOP\textquotesingle s stronger performance is therefore consistent
with its temporal summarisation being better matched to these brief
intervals.

In the regenerated primary benchmark, FlowMOP had higher sensitivity
than FlowCut in both Segment settings and in the 5000-event Bimix and
Trimix benchmarks. At 5000 events, FlowMOP also had higher specificity
than both competitors across Segment, Bimix, and Trimix. In the smaller
mixed-source benchmarks, FlowMOP retained higher specificity than
PeacoQC without a significant sensitivity disadvantage.

Consequently, the earlier smoothing attribution has been removed. The
matched Time-only benchmark directly supports the distinction between
FlowMOP\textquotesingle s fluorescence-population-summary decision and
FlowCut\textquotesingle s response to acquisition-density changes when
fluorescence is unchanged. PeacoQC\textquotesingle s documented bin-size
trade-off provides the relevant context for its fixed-setting comparison
with FlowMOP.

We thank the reviewers for prompting this direct examination of
smoothing. We tested a range of smoothing settings across all synthetic
time-gating inputs while holding all other FlowMOP settings fixed (Table
S4). No smoothing provided the best specificity, but at the cost of
sensitivity. Inspection showed that the additional target-source events
removed under smoothing lay in distribution shoulders, where target and
non-target events were difficult to distinguish reliably. We therefore
retained smoothing as the conservative default: it avoids treating
ambiguous shoulder events as confidently valid and provides higher
sensitivity, while accepting the corresponding specificity trade-off.
Among the smoothed settings, \texttt{0.01,0.05} provided the strongest
balance and was selected as the default. We reran FlowMOP across all
Figure 2 inputs and regenerated Figure 2 using these outputs.

All algorithms were compared using documented recommended, default, or
automatically selected settings, including fixed FlowMOP parameters, to
reflect typical unsupervised use. We did not perform extensive parameter
tuning for FlowCut or PeacoQC because the original method descriptions
do not provide dataset-specific guidance for how such tuning should be
performed. We therefore restrict our conclusions to the fixed comparison
settings and the behaviours directly tested by the matched Time-only
benchmark.

The two smoothing resolutions are now described as complementary spline
fits applied to the same time-bin fluorescence-summary series before MAD
filtering. A bin can be detected through either smoothing pass; the
manuscript no longer relies on the imprecise phrase ``multiple types of
aberrations.''

\subsubsection{Changes made}\label{changes-made-3}

\begin{itemize}
\item
  We removed the earlier smoothing attribution and added the following
  matched comparison:

  \begin{quote}
  ``To test the effect of flow-rate disturbances aligned with or
  independent of source-linked fluorescence changes, we altered only the
  Time channel while leaving fluorescence, scatter, source labels, and
  event order unchanged (Fig. S4). FlowMOP and PeacoQC were unchanged
  under both source-linked and random Time warping. In contrast,
  FlowCut\textquotesingle s sensitivity and specificity shifted after
  Time-only perturbation.''
  \end{quote}
\item
  We clarified the function of the two smoothing resolutions as follows:

  \begin{quote}
  ``Two spline smoothing values, one small and one larger (current
  default \texttt{0.01,0.05}), are applied to the returned time-bin
  series before median absolute deviation (MAD) filtering.''
  \end{quote}

  \begin{quote}
  ``The two smoothing resolutions target both shorter and more sustained
  deviations, while parameter voting limits removal driven by isolated
  noisy channels in higher-dimensional panels.''
  \end{quote}
\item
  We expanded the smoothing analysis across all 173 primary synthetic
  time-gating inputs. Supplementary Table S4 reports the tested range
  and identifies \texttt{0.01,0.05} as the strongest-performing smoothed
  setting on the balanced comparison.
\item
  We reran FlowMOP using the selected \texttt{0.01,0.05} setting and
  regenerated the Figure 2 violin distributions, p-values, and
  significance brackets from these outputs.
\item
  We clarified the scope of the fixed-setting comparison:

  \begin{quote}
  ``The primary comparison used recommended or automatically selected
  settings, including fixed FlowMOP parameters, to reflect typical
  unsupervised use.''
  \end{quote}
\item
  We added the complete mechanism-benchmark methods and results as
  Figure S4.
\end{itemize}

\textbf{Location:} Methods, ``Time-only acquisition-rate mechanism
benchmark'' and ``MAD-smoothing ablation and default selection'';
Results, ``Time Gating'' and ``Synthetic Time Gating Benchmark''; Figure
2; Figure S4; Supplementary Table S4; Discussion, ``Synthetic Sample
Time Gating.''

\subsection{Combined Comment 5 --- Parameter Voting and the
Ten-Parameter
Threshold}\label{combined-comment-5--parameter-voting-and-the-ten-parameter-threshold}

\textbf{Raised by:} Reviewer 2, Comments 18 and 21.

\subsubsection{Reviewer 2, Comment 18}\label{reviewer-2-comment-18}

\begin{quote}
The authors say ''when \textgreater10 parameters are present, bins are
rejected if they are flagged by two or more parameters''. Why ``10'' is
used as threshold here?
\end{quote}

\subsubsection{Reviewer 2, Comment 21}\label{reviewer-2-comment-21}

\begin{quote}
The authors say, ``parameter voting to ensure that higher dimensional
sample sets are not overly degraded.'' It would be nice to show how bad
the overly degradation will be is parameter voting is not applied. It is
a little a bit concern that events with aberrant signal in one
fluorescence are not detected, even in a big panel.
\end{quote}

\subsubsection{Response}\label{response-4}

The revised manuscript identifies the ten-parameter cutoff as an
empirical operational threshold rather than a mathematically derived
boundary.

The rationale for requiring two channel flags in panels with more than
ten eligible parameters is conservative. Biological datasets have a
non-zero probability of isolated channel-specific outliers, and the
opportunity for such outliers increases with panel dimensionality.
Requiring corroboration from a second parameter reduces the risk that a
single noisy channel removes otherwise valid biological events. The
corresponding trade-off is that an abnormality confined to one
fluorescence channel may be retained in a high-dimensional panel. We did
not add a separate voting-ablation analysis and have avoided presenting
the rule as universally optimal.

\subsubsection{Changes made}\label{changes-made-4}

\begin{itemize}
\item
  We added the following rationale and limitation:

  \begin{quote}
  ``For panels with more than 10 parameters, FlowMOP requires two or
  more parameters to flag a bin before rejection. This empirical
  safeguard reduces false-positive removal caused by isolated noisy
  channels in high-dimensional panels, although an aberration confined
  to one channel may consequently be retained (Figure 1A).''
  \end{quote}
\item
  We did not add a voting-ablation experiment and do not present the
  threshold as universally optimal.
\end{itemize}

\textbf{Location:} Results, ``Time Gating.''

\subsection{Combined Comment 6 --- Expert Rankings, Terminology,
Visualisation, and Bayesian
Modelling}\label{combined-comment-6--expert-rankings-terminology-visualisation-and-bayesian-modelling}

\textbf{Raised by:} Associate Editor; Reviewer 1, Comment 3; Reviewer 2,
Comment 16.

\subsubsection{Associate Editor --- Expert Preference Rankings and
Statistical
Interpretation}\label{associate-editor--expert-preference-rankings-and-statistical-interpretation}

\begin{quote}
My other major concern pertains to the substantial space --- both in the
main manuscript and supplemental materials --- devoted to preferential
rankings by four experts across multiple sample types and algorithms,
analyzed using a Bayesian statistical framework. Setting aside the
question of whether all experts can reasonably be considered equally
expert across all sample types, and the limited statistical power
afforded by four data points, I find the interpretation of these
rankings fundamentally ambiguous.

As expected, experts generally preferred human-gated data over automated
tools --- but this finding is difficult to interpret in isolation.
Forcing experts to rank outputs necessarily produces an ordering, but
that ordering conflates two very different scenarios: one in which all
automated results are unacceptable and the other in which all automated
results are perfectly adequate, with human gating simply preferred at
the margin. Similarly, even if FlowMop was ranked above the other
automated tools, the reader has no basis for judging whether this
difference reflects a meaningful improvement in analytical utility or
merely a marginal statistical preference with no practical consequence.

I would suggest that the authors replace or supplement the ranking
analysis with a simpler, more interpretable quality scale --- for
example: 1 = excellent gating; 2 = adequate and suitable for analysis; 3
= inadequate for analysis but not terrible; 4 = unacceptable. This
approach would convey absolute quality rather than relative preference
and would be more informative to readers.
\end{quote}

\subsubsection{Reviewer 1, Comment 3 --- Refinement of Human Subjective
Rankings}\label{reviewer-1-comment-3--refinement-of-human-subjective-rankings}

\begin{quote}
While the subjective rankings provide an interesting window into expert
perspectives on automated gating, the presentation and terminology in
this section require significant refinement:

Nomenclature: In computer science, the term ``benchmarking'' typically
implies an objective, repeatable measure of performance in specific
scenarios. In the context of subjective human preference, ``Comparison''
or ``Evaluation'' is a more accurate/appropriate term. This sentence,
``FlowMOP's outputs in these samples were benchmarked against a set of
human experts'', should more precisely say that
FlowMOP\textquotesingle s outputs were ``compared against expert
gating''.

Methodological Robustness: Utilizing a single expert per tissue type
limits the statistical power of the ranking methodology. It prevents the
reader from seeing the intra-tissue consensus---how consistently
multiple experts on the same tissue agree on the rank-order of methods
used on a specific sample. What it currently highlights is how a single
tissue expert has a unique perspective of the correct gating for a given
tissue and inter-operator gating variability; not algorithmic gating
superiority. Other improvements could be adding different processing
conditions (e.g. fixation) or more variation in tissue used (e.g.
include human PBMC or other species\textquotesingle{} tissues).

Data Visualization: The ranking scale (7 = Best, 1 = Worst) is possibly
counter-intuitive and should be reversed to follow standard ordinal
ranking conventions. The row axis title in the grid figure should be
updated to ``Gate Provided By'' to clearly disambiguate between the
creator of the gate (e.g., ``Expert 1'') and the individual performing
the ranking (`` Expert'').

The scatter plot in Panel B adds limited value; replacing it with an
``Average Score'' column at the end of the grid would be cleaner. This
statement is oddly contorted to be true, but is misconstruing the data:
``However, it is of note that in 4/5 and 7/9 datasets in the debris and
doublets tasks respectively, FlowMOP was not the least preferred,
indicating superiority to at least one human benchmarker (Fig. 6A,
7A)''. Another interpretation is that FlowMOP had the lowest average
score in Fig. 7A/B and had the fewest tissue cases where experts deemed
it as the best result or even better-than-average traditional gating.
\end{quote}

\subsubsection{Reviewer 2, Comment 16}\label{reviewer-2-comment-16}

\begin{quote}
There is only a ``Bayesian Modelling'' section under the ``Non-synthetic
samples'' section in the Method. The Bayesian Modelling section
described the method to rank performance of human experts, FlowMOP,
FlowCut and PeacoQC. It is not about how ``Non-synthetic samples'' are
generated at all. Also, the purpose of the Bayesian Modelling is well
described in the beginning of this subsection, which makes it difficult
to follow.
\end{quote}

\subsubsection{Response}\label{response-5}

Forced rankings measure relative preference rather than objective
event-level accuracy. We therefore replaced ``benchmarking'' with
``expert comparison'' or ``expert evaluation,'' removed the contorted
``not least preferred'' interpretation, and now state directly how
FlowMOP ranked relative to the human and automated gates.

The evaluation produced distinct results across the three cleaning
modules. For time gating, FlowMOP had the best mean rank among the
automated methods and was preferred to FlowCut (BF = 5.39, P = 84.3\%)
and PeacoQC (BF = 12.10, P = 92.4\%). Human gates generally ranked
higher for debris and doublet removal, although FlowMOP was competitive
in several tissue-specific comparisons: it ranked first for debris
removal in mouse blood, third for debris removal in human liver and
mouse skin, and second for doublet removal in human liver (Figs. S5-S7).
We now report these principal findings and statistics in the main
Results; the complete rankings and statistical comparisons remain in the
Supplementary Results and Tables S1B-E.

The revised Discussion considers these results alongside the other
validation approaches. The preference for FlowMOP over the automated
time-gating comparators agrees with its combined sensitivity-specificity
performance in the synthetic benchmarks. The more variable debris and
doublet rankings are discussed in relation to their dependence on
sample-specific scatter distributions. Taken together, the results
indicate that FlowMOP-generated gates can be comparable in acceptability
to human gating across many use cases.

The evaluation was designed as a relative-ranking exercise, with each
dataset assessed by an expert familiar with that sample type. We did not
retrospectively replace this with an absolute-quality scale because
doing so would require repeating the assessment under a different
protocol. Instead, we retain the ranking analysis and report its
outcomes directly, alongside the source-labelled synthetic controls and
the new PBMC and tumour analyses. The reviewer\textquotesingle s request
for broader biological datasets, including human PBMCs, and the
Associate Editor\textquotesingle s request for downstream biological
endpoints are addressed together in Combined Comment 3.

The visual summaries now include an Average Score column within each
ranking grid rather than a separate panel, and the row axis is labelled
``Gate Provided By.'' The displayed scale now follows the standard
convention that rank 1 is best. This display correction did not change
the underlying ranking order or Bayesian analysis.

The Bayesian model is retained because the observed outcome is an
ordinal ranking. The Plackett--Luce model retains the complete ordering
and represents uncertainty in relative preference without treating
differences between adjacent ranks as equal-interval continuous
measurements. We have clarified that it estimates relative preference
only. The subsection is presented as the statistical analysis of the
expert-ranking data rather than as a method for generating non-synthetic
samples.

\subsubsection{Changes made}\label{changes-made-5}

\begin{itemize}
\item
  We defined the scope of the expert-ranking analysis as follows:

  \begin{quote}
  ``The expert-ranking analysis was used to summarize relative
  preferences among gates generated by FlowMOP, comparator algorithms,
  and human operators; it was not treated as an absolute measure of
  gating adequacy.''
  \end{quote}
\item
  We revised the Supplementary Results terminology and interpretation as
  follows:

  \begin{quote}
  ``FlowMOP's outputs in these samples were compared with
  expert-provided gates using a forced-ranking preference task. These
  rankings measure relative expert preference and should not be
  interpreted as an absolute measure of gating adequacy.''
  \end{quote}
\item
  We retained the Plackett--Luce analysis and clarified that it models
  ordinal preference rankings rather than absolute gating quality.
\item
  We moved the response concerning dataset breadth and downstream
  biological validation to Combined Comment 3 and cross-reference it
  here rather than treating the expert rankings as biological
  validation.
\item
  We retained the full expert-ranking analysis in the Supplementary
  Information and added its principal comparative findings and
  statistics to the main Results, with their implications discussed
  alongside the synthetic and biological-validation findings.
\item
  We regenerated Supplementary Figures S5--S7, changed the row-axis
  title to ``Gate Provided By,'' displayed rank 1 as best, and replaced
  the separate summary panel with an Average Score column in each grid.
\item
  We revised the Figure S5 caption as follows:

  \begin{quote}
  ``Figure S5. Expert preference rankings for time gates provided by
  four human experts, FlowMOP (black border), FlowCut, and PeacoQC
  across nine datasets. Rows indicate the gate provider, and the final
  column shows the mean rank across datasets. Rank 1 indicates greatest
  preference; therefore, a lower average score indicates greater
  preference. Abbreviations: DRG, dorsal root ganglion; CNS, central
  nervous system.''
  \end{quote}
\item
  We clarified the scoring direction in the Supplementary Results and
  figure captions as follows:

  \begin{quote}
  ``Rank 1 indicates the greatest preference.''
  \end{quote}
\end{itemize}

\textbf{Location:} Results, ``Expert preference evaluation'';
Discussion; Conclusion; Supplementary Information, ``Supplementary
expert preference evaluation,'' including ``Supplementary methods'' and
``Supplementary results''; Supplementary Figures S5--S7; Supplementary
Tables S1B-E.

\subsection{Combined Comment 7 --- Dataset Scale and
Complexity}\label{combined-comment-7--dataset-scale-and-complexity}

\textbf{Raised by:} Associate Editor.

\subsubsection{Associate Editor --- Smaller
Point}\label{associate-editor--smaller-point}

\begin{quote}
In the Introduction, the authors state that sophisticated analysis
software has contributed to datasets of increasing size and complexity.
It is not clear to me how analysis software makes debris or time gating
more complex, particularly when pre-processing is predominantly
scatter-based. Do the authors perhaps mean that derived parameters from
imaging cytometry increase pre-processing complexity? Similarly, the
claim that ``datasets exceed human capabilities'' presumably refers to
the sheer number of files rather than the complexity of individual
datasets --- this should be stated explicitly. It is also not apparent
from the manuscript how the tools described here depend on data
complexity per se.
\end{quote}

\subsubsection{Response}\label{response-6}

The revised Introduction distinguishes study scale from the intrinsic
complexity of an individual preprocessing gate. It now refers explicitly
to increasing numbers of files, events, and measured parameters and
explains that this scale makes repeated manual preprocessing
time-consuming and often impractical. It no longer implies that analysis
software itself necessarily makes a debris or time gate more complex.

\subsubsection{Changes made}\label{changes-made-6}

\begin{itemize}
\item
  We clarified the distinction between study scale and individual-gate
  complexity as follows:

  \begin{quote}
  ``Modern flow-cytometry studies can contain increasing numbers of
  files, events, and measured parameters {[}2{]}. The large scale of
  these data makes repeated manual preprocessing time-consuming and
  often impractical, potentially amplifying variability between
  operators. Reproducible automated preprocessing is therefore valuable
  for large studies.''
  \end{quote}
\end{itemize}

\textbf{Location:} Introduction.

\subsection{Combined Comment 8 --- Acquisition Settings and Newer
Scatter
Configurations}\label{combined-comment-8--acquisition-settings-and-newer-scatter-configurations}

\textbf{Raised by:} Reviewer 2, Comments 13 and 28.

\subsubsection{Reviewer 2, Comment 13}\label{reviewer-2-comment-13}

\begin{quote}
The voltage set of the cytometer significantly impact the resolution of
debris. If the voltage is relatively low with bad resolution, can
FlowMOP perform the automatic task well? Also, please describe the
voltage setting in the data generation section in the Method.
\end{quote}

\subsubsection{Reviewer 2, Comment 28}\label{reviewer-2-comment-28}

\begin{quote}
The Xenith flow cytometry form Thermo Fisher provides 405 FSC/SSC, 488
FSC/SSC, and polar 488 FSC/SSC. This provides more flexible approach to
gate target population. How can the FlowMOP adapt to this kind of flow
data? It would be nice to discuss this in the Discussion.
\end{quote}

\subsubsection{Response}\label{response-7}

The benchmarking relied on real, existing datasets, each acquired using
experiment-specific, biology-driven antibody panels and instrument
settings. Consequently, acquisition voltage/gain settings differed among
datasets and were not controlled variables in this study. As with the
antigen--fluorochrome combinations, these settings are therefore not
material to the preprocessing comparisons reported here, and a single
voltage/gain specification would not describe the datasets used. We have
clarified this scope in the Methods.

Acquisition settings nevertheless determine whether the scatter signals
needed by the debris and doublet modules are adequately resolved.
FlowMOP requires appropriate acquisition voltage/gain settings; if
relevant signals are poorly resolved or saturated, the lost information
cannot be recovered and reliable cleaning cannot be guaranteed. FlowMOP
currently expects users to identify the scatter channels used by the
workflow; it has not been validated systematically across 405-nm,
488-nm, and polar 488-nm FSC/SSC configurations on instruments with
multiple scatter measurements. Future versions could assess multiple
scatter-channel pairs and select or combine the pair with the clearest
debris/doublet separation.

\subsubsection{Changes made}\label{changes-made-7}

\begin{itemize}
\item
  We clarified the dataset-specific acquisition settings in the Methods:

  \begin{quote}
  ``The benchmarking uses real, existing datasets, each acquired using
  experiment-specific, biology-driven antibody panels and instrument
  settings. Consequently, the antigen--fluorochrome combinations and
  acquisition voltage/gain settings differ among datasets and were not
  variables under evaluation; they are therefore not material to the
  preprocessing comparisons addressed here.''
  \end{quote}
\item
  We retained the following combined acquisition/scatter limitation:

  \begin{quote}
  ``FlowMOP requires appropriate acquisition voltage/gain settings; if
  relevant signals are poorly resolved or saturated, the lost
  information cannot be recovered and reliable cleaning cannot be
  guaranteed. FlowMOP currently expects users to identify the scatter
  channels used by the workflow; it has not been validated
  systematically across 405-nm, 488-nm, and polar 488-nm FSC/SSC
  configurations on instruments with multiple scatter measurements.
  Future versions could assess multiple scatter-channel pairs and select
  or combine the pair with the clearest debris/doublet separation.''
  \end{quote}
\end{itemize}

\textbf{Location:} Methods, ``Preparation of Human PBMC Samples for
Synthetic Time Benchmarking Samples''; Discussion, ``Other remarks.''

\subsection{Combined Comment 9 --- Temporal Artifacts and Synthetic
Time-Sample
Design}\label{combined-comment-9--temporal-artifacts-and-synthetic-time-sample-design}

\textbf{Raised by:} Associate Editor; Reviewer 2, Comments 3, 4, and 14.

\subsubsection{Associate Editor --- Smaller
Point}\label{associate-editor--smaller-point-1}

\begin{quote}
Regarding the bimix and trimix data used to simulate realistic
fluorescence fluctuations: I am uncertain how these mixtures reproduce
the kinds of flow rate variations that occur during actual experiments.
If I understand correctly, the simulated fluctuations were subtle enough
that human reviewers could not reliably detect them. I would ask the
authors to clarify whether this represents a clinically or
experimentally common problem --- in my experience, flow rate deviations
typically manifest as clear departures at the beginning or end of an
acquisition, with mid-acquisition instabilities being relatively
uncommon.
\end{quote}

\subsubsection{Reviewer 2, Comment 3}\label{reviewer-2-comment-3}

\begin{quote}
The authors should explain the concept of ``temporal artifacts'' in the
abstract, given this is the core issue to be addressed in the study.
\end{quote}

\subsubsection{Reviewer 2, Comment 4}\label{reviewer-2-comment-4}

\begin{quote}
There are many terms in the abstract ``(concatenated and mixed time
perturbations, high-debris mixtures, and CFSE/CTV co-labeled
doublets)''. I am not sure whether it is a good idea to put them here
without explanation, considering unexplained term might add confusion to
the readers. Anyway, the authors should introduce these terms in the
introduction, which is missing.
\end{quote}

\subsubsection{Reviewer 2, Comment 14}\label{reviewer-2-comment-14}

\begin{quote}
The ``Generation of Synthetic Time Samples'' section is difficult to
understand. This section is important for readers to understand the
validation results. I recommend the authors use a illustration to show
these 3 synthetic manners. Also, please explain why these three manners
are designed in this way? So that the readers can follow the text.
\end{quote}

\subsubsection{Response}\label{response-8}

The revised text explains the practical relevance of the Bimix and
Trimix designs. It defines a ``microblockage'' operationally as a short,
self-resolving mid-acquisition disturbance that produces a localized
fluorescence shift; the term does not imply that a physical obstruction
was directly observed. Existing work supports the substance of this
phenomenon even though it uses broader terminology. flowAI documented
flow-rate surges interspersed throughout acquisitions and their
association with signal-intensity variation {[}10{]}. PeacoQC describes
temporary acquisition shifts caused by slow sample uptake or a clog and
notes that such problems can be difficult to detect manually {[}5{]},
while FlowCut describes manual identification and removal of transient
acquisition problems as time-consuming and subjective {[}9{]}.

Because the affected intervals can be short and their events
interspersed with otherwise plausible events, they are difficult to
recognize reliably by eye and impractical to exclude using a series of
manual gates. We use ``microblockage'' for this subtle, self-resolving
acquisition failure mode that a time-quality-control method should be
able to detect. Estimating its prevalence across experiments would
require a separate study.

The Bimix and Trimix samples model the observable fluorescence
consequence in this operational definition rather than reproducing or
proving a physical obstruction. This distinction is important. The
synthetic samples can appear acceptable by eye, but their retained
source labels reveal short intervals containing events from an
intentionally perturbed fluorescence source. Under the benchmark
definition, those events should be excluded even when visual inspection
alone would not identify a defensible manual gate. The apparent visual
normality of these files is therefore the reason an event-labelled
benchmark is needed.

The downstream biological-validation analyses are presented together in
Combined Comment 3. That evidence provides an independent biological
context for the technical benchmark; it is cross-referenced here rather
than repeated. Together, the synthetic and biological analyses motivate
an evaluation that includes both readily visible sustained artifacts and
visually subtle transient microblockages.

We also defined temporal artifacts, simplified the terminology, and
added Figure S1. We clarified that the primary synthetic samples model
fluorescence changes across acquisition order without flow-rate
disturbances, because rate changes without corresponding fluorescence
changes should not prompt event exclusion. Flow-rate effects were tested
separately, either aligned with source-linked fluorescence changes or
independently of them. FlowMOP was unchanged in both conditions, whereas
FlowCut\textquotesingle s removal behavior shifted (Figure S4).

\subsubsection{Changes made}\label{changes-made-8}

\begin{itemize}
\item
  In the Abstract, we incorporated the definition into the
  methodological summary:

  \begin{quote}
  ``Methodologically, FlowMOP identifies temporal
  artifacts---acquisition-dependent deviations in event quality or
  fluorescence signal---via parameter-wise peak checks, bin-level
  fluorescence summaries across acquisition time, and robust outlier
  rejection.''
  \end{quote}
\item
  In the Methods, we added the following clarification:

  \begin{quote}
  Here, we use ``microblockage'' operationally to denote a short,
  self-resolving mid-acquisition disturbance that produces a localized
  fluorescence shift; the term does not imply that a physical
  obstruction was directly observed. Segmented samples model sustained
  changes, whereas Bimix and Trimix model the observable fluorescence
  consequence in this operational definition by introducing short
  source-defined fluorescence shifts during acquisition. They do not
  recreate or establish the physical mechanism itself. The Bimix and
  Trimix files can appear acceptable on visual inspection because the
  altered intervals are short and interspersed with otherwise plausible
  events; however, the retained source labels identify the intentionally
  perturbed events that should be excluded under the benchmark
  definition.
  \end{quote}

  \begin{quote}
  ``Flow-rate disturbances without corresponding fluorescence changes
  should not prompt event exclusion; these samples therefore model
  fluorescence changes across acquisition order without introducing
  flow-rate disturbances. Flow-rate effects were tested separately by
  altering Time either in alignment with source-linked fluorescence
  changes or independently of them.''
  \end{quote}
\item
  In the Results, we added the following comparison:

  \begin{quote}
  ``To test the effect of flow-rate disturbances aligned with or
  independent of source-linked fluorescence changes, we altered only the
  Time channel while leaving fluorescence, scatter, source labels, and
  event order unchanged (Fig. S4). FlowMOP and PeacoQC were unchanged
  under both source-linked and random Time warping. In contrast,
  FlowCut\textquotesingle s sensitivity and specificity shifted after
  Time-only perturbation.''
  \end{quote}
\item
  In the Discussion, we now explain why the visually subtle synthetic
  cases are experimentally relevant and require an event-labelled
  benchmark:

  \begin{quote}
  Transient acquisition disturbances are not confined to the beginning
  or end of a run: flow-rate surges interspersed throughout acquisition
  and associated signal-intensity variation have been documented
  {[}10{]}. PeacoQC notes that temporary acquisition problems can be
  difficult to detect manually {[}5{]}, and FlowCut describes manual
  identification and removal of transient acquisition problems as
  time-consuming and subjective {[}9{]}. We use ``microblockage''
  operationally for a short, self-resolving instance of this broader
  phenomenon that produces a localized fluorescence shift, without
  asserting that a physical obstruction was directly observed.
  \end{quote}

  \begin{quote}
  ``Although these synthetic samples can appear acceptable on visual
  inspection, the source labels show that the short altered intervals
  contain events from an intentionally perturbed fluorescence source and
  therefore should be excluded under the benchmark definition. Visual
  subtlety is thus a central feature of the benchmark: it demonstrates
  why apparent normality by eye is not sufficient ground truth.''
  \end{quote}
\item
  We cross-reference the consolidated downstream biological-validation
  response in Combined Comment 3 rather than repeating its analyses
  here.
\item
  We added Figure S1 with the following caption:

  \begin{quote}
  ``Figure S1: Construction of Segment, Bimix, and Trimix synthetic time
  samples. No flow-rate disturbance was introduced.''
  \end{quote}
\end{itemize}

\textbf{Location:} Abstract; Introduction; Methods, ``Generation of
Synthetic Time Samples'' and ``Time-only acquisition-rate mechanism
benchmark''; Results, ``Synthetic Time Gating Benchmark''; Discussion,
``Synthetic Sample Time Gating''; Figures S1 and S4. The
biological-validation changes are detailed in Combined Comment 3.

\subsection{Combined Comment 10 --- Debris Methodology and
Validation}\label{combined-comment-10--debris-methodology-and-validation}

\textbf{Raised by:} Reviewer 1, Comment 2; Reviewer 2, Comments 15, 22,
and 24.

\subsubsection{Reviewer 1, Comment 2 --- Debris Gating
Methodology}\label{reviewer-1-comment-2--debris-gating-methodology}

\begin{quote}
The inclusion of an automated debris-gating method is a welcome
addition, particularly as this is often omitted in other QC tools.
However, the current reliance on FSC-A as the primary determinant for
debris may explain the lower rankings in expert consensus (Figure 6A).

Critique: While the rationale is that debris is generally ``low-size,''
this does not account for larger morphological debris, such as
large-cell clumps (e.g. tumor) or non-biological aggregates, which can
occupy higher FSC/SSC ranges. Metadata-based margin removal techniques
may augment this debris filtering.

Recommendation: The authors should provide a more robust rationale for
leaning exclusively on FSC-A or, ideally, discuss how the integration of
SSC-A or Pulse Width might improve debris discrimination beyond
ultra-low-size events.
\end{quote}

\subsubsection{Reviewer 2, Comment 15}\label{reviewer-2-comment-15}

\begin{quote}
The authors say ``For assessment of FlowMOP gating performance, `high
debris' and `low debris' samples were synthetically combined.'' How are
they combined? Please remember, well-described Method guarantees the
transparency of the study.
\end{quote}

\subsubsection{Reviewer 2, Comment 22}\label{reviewer-2-comment-22}

\begin{quote}
The FlowMOP applies an FSC-A based threshold to exclude debris. I have
concern about this FSC-A based idea. Commonly, we use FSC vs. SSC
scatter to gate debris, which is clear in samples like PBMC. However,
the resolution of a simple FSC-A histogram will not be as good as that
in a FSC vs. SSC scatter, because debris and cells will overlap in FSC-A
axis. In the Figure 3, the authors show representative results which is
based on simple samples, which is possible to use the FSC-A based
threshold method to gate debris. However, in difficult cases, which is
not very rare, this FSC-A based threshold method may perform very bad.
\end{quote}

\subsubsection{Reviewer 2, Comment 24}\label{reviewer-2-comment-24}

\begin{quote}
I would like to see the debris remove rate calculated against the human
expert.
\end{quote}

\subsubsection{Response}\label{response-9}

The revised Methods describes the debris-removal procedure more
precisely. FlowMOP ultimately applies a one-dimensional FSC-A threshold,
but that threshold is not obtained from a single unconditioned FSC-A
histogram. The algorithm identifies eligible fluorescence parameters,
examines the FSC-A structure of each parameter\textquotesingle s
positive population, derives a candidate FSC-A threshold from each
eligible channel, and applies the median candidate threshold to the
sample. Thus, information across eligible fluorescence channels informs
the final FSC-A gate, while the applied decision remains a conservative
FSC-A threshold.

The current module is not a universal debris classifier. Debris and
intact cells can overlap on FSC-A, and large clumps, tumour-associated
material, non-biological aggregates, or other high-scatter contaminants
may not be removed reliably. We specifically considered whether SSC-A
should be incorporated. SSC-A can be informative for large or internally
complex debris, but large and internally complex events can also
represent desirable populations that should be retained. Their
interpretation varies with tissue, staining panel, cytometer
configuration, acquisition settings, and the biological populations of
interest. Unlike the comparatively transferable relationship between
very low FSC-A and small debris, there is no broadly safe high-SSC-A
rule for excluding large events.

In the present technical-control datasets, the labelled small-debris and
desirable source populations overlapped substantially along SSC-A, so a
fixed SSC-A threshold offered limited additional separation for the
small-event removal targeted by FlowMOP. We therefore did not
incorporate SSC-A into the current debris decision. Pulse-width
measurements may similarly help identify some aggregates or clumps, but
their availability and interpretation are instrument- and
acquisition-dependent. Incorporating SSC-A or pulse width responsibly
would require a wide configurable parameter range or sample-specific
multivariate decision rules validated for the intended dataset, panel,
instrument, and cell populations. Metadata-based margin removal is
complementary rather than a substitute for debris classification: it can
identify events at acquisition limits, but it does not identify all
low-FSC debris or large aggregates. FlowMOP currently includes only an
FSC-A maximum-value precleaning check and does not implement a
generalized margin-event filter. The revised Discussion therefore
defines the current scope as conservative low-FSC debris removal and
presents SSC-A, pulse-width measurements, margin-event metadata, and
sample-specific multivariate models as future extensions rather than
universal default filters.

In debris-poor sample 6A, we discovered that the FSC-A peak-based
procedure was liable to misclassify a lower-FSC biological population
composed predominantly of Live CD45+ T cells as debris. This illustrated
a limitation of the procedure: when a sample lacks a debris population,
the module may remove biologically relevant events and its output should
be inspected. The complete downstream results are reported in the
revised Results, Figure 6, and Combined Comment 3.

The synthetic preparation is now described explicitly. Approximately
equal event numbers were sampled from the high-debris and low-debris
sources and concatenated into matched mixtures while retaining source
labels. These source labels provide the objective basis for the reported
post-cleaning proportions.

The requested human comparison is represented by the percentage-based
analysis in Figure 3. FlowMOP\textquotesingle s enrichment of the
labelled low-debris component is compared directly with four expert
gates, with mean percentages, variability, and paired statistical
comparisons. Because the input event count differs between samples,
percentage-based retention and enrichment are more directly comparable
than raw removed-event counts.

FlowMOP determines its debris gate independently for each sample,
whereas manual debris gating is commonly performed groupwise by applying
one gate across related samples. We therefore also asked the experts to
perform both groupwise and individual-sample gating. No difference was
detected between these strategies for any expert (Fig. 3D, unadjusted
paired t-tests, p \textgreater{} 0.05), indicating that the comparison
with FlowMOP was not driven by this difference in gating strategy.

\subsubsection{Changes made}\label{changes-made-9}

\begin{itemize}
\item
  We described the synthetic debris mixture as follows:

  \begin{quote}
  ``For assessment of FlowMOP debris-gating performance, high-debris and
  low-debris samples were combined by sampling approximately equal event
  numbers from each source and concatenating them into matched synthetic
  mixtures while retaining source labels for ground-truth
  quantification.''
  \end{quote}
\item
  We clarified the scope and derivation of the FSC-A gate as follows:

  \begin{quote}
  ``FlowMOP's debris module targets small, low-FSC debris. The final
  gate is applied on FSC-A, but its threshold is informed by FSC-A
  distributions across eligible fluorescence-positive populations rather
  than the overall FSC-A histogram alone.''
  \end{quote}

  \begin{quote}
  ``SSC-A may improve recognition of larger or internally complex
  debris, while pulse-width measurements may assist with aggregates.
  Accordingly, FlowMOP does not currently incorporate SSC-A or
  pulse-width measurements into its debris decision because broadly
  applicable thresholds for these features are difficult to establish
  across tissues, panels, instruments, and acquisition settings.''
  \end{quote}

  \begin{quote}
  ``The median FSC-A threshold across all parameters is taken as the
  final FSC-A gate to be applied to the sample (Figure 1B).''
  \end{quote}
\item
  We expanded the Discussion of SSC-A and pulse-width integration:

  \begin{quote}
  ``SSC-A may improve recognition of larger or internally complex
  debris, while pulse-width measurements may assist with aggregates.
  Accordingly, FlowMOP does not currently incorporate these features
  because broadly applicable decision rules are difficult to establish
  across tissues, panels, instruments, and acquisition settings.''
  \end{quote}
\item
  We clarified that generalized metadata-based margin removal is not
  currently implemented beyond the FSC-A maximum-value precleaning check
  and is a potential complementary extension rather than a complete
  debris classifier.
\item
  We added the debris-poor PBMC result as a direct limitation of the
  current FSC-A peak assumption and contrast it with successful removal
  in the source-labelled benchmark when genuine separable debris is
  present.
\item
  We reported the existing expert comparison as follows:

  \begin{quote}
  ``FlowMOP did not differ significantly from any human evaluator except
  Expert 4, for whom FlowMOP removed more labelled high-debris events
  (Bonferroni-adjusted paired t-test, p = 0.04) (Fig. 3C).''
  \end{quote}
\item
  We clarified that FlowMOP estimates debris gates independently for
  each sample and reported the existing
  groupwise-versus-individual-sample control, for which no difference
  was detected for any expert (Fig. 3D).
\end{itemize}

\textbf{Location:} Methods, ``Synthetic Debris Sample Preparation and
Generation'' and ``Biological-validation analysis''; Results, ``Debris
Gating,'' ``Synthetic Debris Gating Benchmark,'' and ``Biological
validation''; Discussion, ``Synthetic Sample Debris and Doublet Gating''
and ``Biological validation''; Figures 3 and 6.

\subsection{Combined Comment 11 --- Aggregates, Doublets, Saturation,
and
Validation}\label{combined-comment-11--aggregates-doublets-saturation-and-validation}

\textbf{Raised by:} Reviewer 2, Comments 5, 25, and 26.

\subsubsection{Reviewer 2, Comment 5}\label{reviewer-2-comment-5}

\begin{quote}
Preprocessing sometimes include aggregates removal. Is it possible for
FlowMOP to handle this task?
\end{quote}

\subsubsection{Reviewer 2, Comment 25}\label{reviewer-2-comment-25}

\begin{quote}
The FlowMOP creates a histogram of the FSC-A/FSC-H ratio to gate
doublets. I also have a concern about this method. One unmentioned
hypothesis of this algorithm is that all cells fall within the
measurement range of FSC. However, it is very common that some cells
exceed the FCS limit and collapsed in the edge of a FSC-A vs FSC-H
scatter plot. For example, big myeloids in PBMC, if the target
population is lymphocytes. Such population is expected to lead to weird
histogram of the FSC-A/FSC-H ratio. This can be difficult to be
addressed by the FlowMOP and are not mentioned in the study.
\end{quote}

\subsubsection{Reviewer 2, Comment 26}\label{reviewer-2-comment-26}

\begin{quote}
The authors use CTV-CFSE double positive cells to represent doublets.
However, there are probably CTV- CTV doublets and CFSE - CFSE doublets,
which are missed in the performance evaluation. I would like to see the
doublet remove rate calculated against the human expert.
\end{quote}

\subsubsection{Response}\label{response-10}

FlowMOP\textquotesingle s doublet module targets events whose
FSC-A/FSC-H and, where available, SSC-A/SSC-H pulse ratios distinguish
them from singlets. It can remove some higher-order coincident events
when those events produce separable ratio structure, as observed for
triplet-like events in the technical controls, but it cannot identify
every aggregate or clump from these ratios alone.

The general requirement for appropriate acquisition settings is
addressed in Combined Comment 8. For the doublet module specifically,
FSC-A/FSC-H and SSC-A/SSC-H ratios must remain informative. Saturated or
edge-collapsed scatter measurements may therefore require acquisition
review, manual intervention, or alternative pulse-shape features. We did
not add a new saturation-handling algorithm.

The reviewer is also correct that CTV-CFSE double-positive events
represent only heterologous labelled doublets. CTV-CTV and CFSE-CFSE
doublets cannot be distinguished from their respective single-labelled
populations using the dye labels alone. The technical-control design
therefore gives a sensitive measure of retained known heterologous
doublets but cannot establish the complete false-positive and
false-negative rate for all doublet classes.

Figure 4 already compares the percentage of CTV-CFSE double-positive
events removed by FlowMOP with the corresponding percentages removed by
human experts. We now describe this endpoint and its same-label
limitation explicitly rather than implying complete doublet ground
truth.

FlowMOP also estimates its doublet gates independently for each sample
rather than applying a shared gate across a group. The experts therefore
performed both groupwise and individual-sample doublet gating. These
strategies did not differ for three of the four experts; Expert 3
removed fewer doublets with individual-sample gating (Fig. 4C, paired
t-test, p = 0.009). We now explain this methodological distinction and
its practical implications in the Results and Discussion.

\subsubsection{Changes made}\label{changes-made-10}

\begin{itemize}
\item
  We cross-reference the general acquisition-setting limitation in
  Combined Comment 8 and retain only the doublet-specific requirement
  that the relevant scatter ratios remain informative.
\item
  We clarified the limits of the CTV/CFSE validation as follows:

  \begin{quote}
  ``Events positive for both CFSE and CTV provide an observable class of
  heterologous labelled doublets, subject to rare dye-transfer events;
  same-label CTV-CTV and CFSE-CFSE doublets are not identifiable from
  these labels alone.''
  \end{quote}
\item
  We reported the comparison with human experts as follows:

  \begin{quote}
  ``No statistically significant difference was detected between FlowMOP
  and any expert for this endpoint (Fig. 4B, paired t-test; unadjusted p
  \textgreater{} 0.05).''
  \end{quote}
\item
  We clarified that FlowMOP estimates doublet gates independently for
  each sample and reported the existing comparison of groupwise and
  individual-sample expert gating (Fig. 4C).
\item
  We retained the observation that FlowMOP removed a triplet-like
  population in the technical-control samples and restricted that
  finding to those samples.
\end{itemize}

\textbf{Location:} Results, ``Doublet Gating'' and ``Synthetic Doublet
Gating Benchmark''; Discussion, ``Synthetic Sample Debris and Doublet
Gating'' and ``Other remarks''; Figure 4.

\subsection{Combined Comment 12 --- CTV/CFSE Preparation
Protocol}\label{combined-comment-12--ctvcfse-preparation-protocol}

\textbf{Raised by:} Reviewer 2, Comment 27.

\subsubsection{Reviewer 2, Comment 27}\label{reviewer-2-comment-27}

\begin{quote}
The protocol about how the CTV-CFSE samples are prepared is not clear.
Are there wash steps in between? This is very important. It would be
nice if the authors can provide detailed protocol as supplementary
files.
\end{quote}

\subsubsection{Response}\label{response-11}

We agree that the original description did not provide enough
information to reproduce the CTV/CFSE preparation, particularly the wash
and quenching steps. We have replaced the abbreviated description with a
detailed protocol specifying the digestion conditions, filtration and
RBC lysis, cell numbers, centrifugation and wash steps, dye stocks and
volumes, labelling and quenching conditions, 1:1 recombination of the
labelled populations, viability staining, and acquisition platform.

\subsubsection{Changes made}\label{changes-made-11}

The revised Methods now state:

\begin{quote}
``To generate samples with high proportions of doublets, C57BL/6 mouse
spleens were injected with 1 mL digestion buffer comprising RPMI
supplemented with 50 µg/mL collagenase P (Roche) and 10 µg/mL DNase I
(Roche). Spleens were incubated for 20 minutes at room temperature (RT)
in a further 1 mL of digestion buffer, mechanically dissociated, and
incubated for a further 20 minutes at RT. Samples were passed through a
70-µm cell strainer with 10 mL FACS wash (PBS containing 2\%
heat-inactivated fetal bovine serum {[}FBS{]}), centrifuged (5 minutes,
500 × \emph{g}, RT), and the supernatant discarded. Cell pellets were
resuspended in 3 mL RBC lysis buffer and incubated for 3 minutes at RT,
after which cells were washed twice with FACS wash and recovered by
centrifugation.

For dye labelling, 5 × 10\^{}6 cells were transferred to each fresh
15-mL tube and stained with either CellTrace Violet (CTV; Invitrogen) or
carboxyfluorescein succinimidyl ester (CFSE; eBioscience). Cells were
centrifuged (5 minutes, 500 × \emph{g}, 4°C), the supernatant was
removed, and pellets were resuspended in 1 mL complete IMDM (cIMDM;
Gibco IMDM supplemented with 10\% heat-inactivated FBS, 100 U/mL
penicillin, 100 µg/mL streptomycin, 2 mM L-glutamine, and 55 µM
2-mercaptoethanol). Two microlitres of 5 mM CTV or 10 mM CFSE were added
to the side of the corresponding tube; samples were rapidly inverted and
mixed, then incubated for 10 minutes at 37°C. Labelling was quenched by
adding 5 mL ice-cold cIMDM and incubating for a further 5 minutes at RT.
Cells were centrifuged and resuspended in cIMDM, after which 2 × 10\^{}6
CTV-labelled cells and 2 × 10\^{}6 CFSE-labelled cells were combined in
a single sample and incubated for 30 minutes at 37°C and 5\% CO2.
Samples were centrifuged, the supernatant was removed, and cells were
stained for 20 minutes at RT with Fixable Viability Dye eFluor 780
(Invitrogen; 1:1,000 in PBS). Cells were washed with PBS, centrifuged,
resuspended in FACS wash, and acquired on a BD LSRII flow cytometer.''
\end{quote}

\textbf{Location:} Methods, ``Synthetic Doublet Sample Preparation and
Generation''.

\subsection{Combined Comment 13 --- Missing and Duplicated Methods
Information}\label{combined-comment-13--missing-and-duplicated-methods-information}

\textbf{Raised by:} Reviewer 2, Comments 11 and 12.

\subsubsection{Reviewer 2, Comment 11}\label{reviewer-2-comment-11}

\begin{quote}
There is duplicated line in the Method ``Samples were acquired on a
Cytek Northern Lights 3-laser (V/B/R) spectral flow cytometer.''
\end{quote}

\subsubsection{Reviewer 2, Comment 12}\label{reviewer-2-comment-12}

\begin{quote}
In the ``Preparation of Human PBMC Samples for Synthetic Time
Benchmarking Samples'' section, the authors mentioned there are many
antibodies in the staining cocktail. However, they did not list which
antibodies.
\end{quote}

\subsubsection{Response}\label{response-12}

The duplicated cytometer sentence has been removed. The benchmarking
uses real, existing datasets acquired with experiment-specific,
biology-driven antibody panels and instrument settings. Consequently,
the antigen--fluorochrome combinations differ among datasets and are not
variables under evaluation; they are not material to the preprocessing
comparisons raised here. We have clarified this scope in the Methods
rather than presenting a single fluorochrome list that would not
describe the distinct panels used across the benchmark datasets.

\subsubsection{Changes made}\label{changes-made-12}

\begin{itemize}
\item
  We removed the duplicated sentence, ``Samples were acquired on a Cytek
  Northern Lights 3-laser (V/B/R) spectral flow cytometer.''
\item
  We clarified why a single fluorochrome list is not applicable across
  the benchmark datasets:

  \begin{quote}
  ``The benchmarking uses real, existing datasets, each acquired using
  experiment-specific, biology-driven antibody panels and instrument
  settings. Consequently, the antigen--fluorochrome combinations and
  acquisition voltage/gain settings differ among datasets and were not
  variables under evaluation; they are therefore not material to the
  preprocessing comparisons addressed here.''
  \end{quote}
\end{itemize}

\textbf{Location:} Methods, ``Preparation of Human PBMC Samples for
Synthetic Time Benchmarking Samples.''

\subsection{Combined Comment 14 --- Precleaning Event Removal and
Threshold}\label{combined-comment-14--precleaning-event-removal-and-threshold}

\textbf{Raised by:} Reviewer 2, Comment 17.

\subsubsection{Reviewer 2, Comment 17}\label{reviewer-2-comment-17}

\begin{quote}
In the ``Precleaning'' section of Results, how many events are being
removed in these used datasets? This information can be provided as
supplementary table/figure to justify the set of threshold (5\%).
\end{quote}

\subsubsection{Response}\label{response-13}

The current implementation uses a 1\% threshold, not 5\%; the
inconsistent earlier value has been corrected. FlowMOP checks the number
of events at the maximum FSC-A value and removes those maximum-valued
events only when they exceed 1\% of the sample.

The number and percentage of events removed by this specific precleaning
step are the relevant quantities. A dataset-level summary of those
values is not yet complete and remains to be added before submission;
percentage-based retention results from later gating stages are not
presented as a substitute for this precleaning summary.

The 1\% cutoff is a pragmatic safeguard rather than a uniquely
ground-truth-derived boundary. It is low enough to identify files with a
non-trivial accumulation of events at the acquisition maximum while
avoiding removal triggered by isolated maximum-valued events. The
revised Results state this operational rationale.

\subsubsection{Changes made}\label{changes-made-13}

\begin{itemize}
\item
  We corrected and defined the threshold as follows:

  \begin{quote}
  ``FlowMOP first checks the input file for events at the limit of
  detection, defined here as events at the maximum FSC-A value for that
  sample. If the number of events at this maximum exceeds a threshold
  (default 1\%), FlowMOP removes these maximum-valued events. Otherwise,
  it retains all values.''
  \end{quote}
\item
  We added the following operational rationale:

  \begin{quote}
  ``The 1\% cutoff is a pragmatic safeguard intended to identify
  non-trivial accumulation at the acquisition maximum without responding
  to isolated maximum-valued events.''
  \end{quote}
\item
  The dataset-level precleaning removal counts and percentages remain to
  be added before submission.
\end{itemize}

\textbf{Location:} Results, ``Precleaning.''

\subsection{Combined Comment 15 --- Figure 1 Axes and
Annotations}\label{combined-comment-15--figure-1-axes-and-annotations}

\textbf{Raised by:} Reviewer 2, Comments 19 and 23.

\subsubsection{Reviewer 2, Comment 19}\label{reviewer-2-comment-19}

\begin{quote}
The Figure 1A missing many axis titles. Also, there are tow subplots
after time-binned median fluorescence. What are these two subplots are
not clearly annotated.
\end{quote}

\subsubsection{Reviewer 2, Comment 23}\label{reviewer-2-comment-23}

\begin{quote}
The Figure 1B also misses many axis titles. The Figure 1C misses y-axis
title.
\end{quote}

\subsubsection{Response}\label{response-14}

Figure 1 remains a conceptual schematic, and its axes are not treated as
quantitative measurements. The revised legend explains the two smoothing
resolutions in Figure 1A, identifies the debris and doublet-ratio
quantities in Figures 1B and 1C, and states that the schematic
fluorescence-intensity and signal-strength axes use arbitrary units.

\subsubsection{Changes made}\label{changes-made-14}

\begin{itemize}
\item
  We revised the Figure 1 legend as follows:

  \begin{quote}
  ``Figure 1. Conceptual schematics depicting FlowMOP's time-gating (A),
  debris-gating (B), and doublet-gating (C) methods; plotted
  fluorescence intensity and signal-strength axes are schematic and use
  arbitrary units. A) FlowMOP selects valid parameters using
  positive-peak detection, generates time-binned fluorescence summaries,
  applies two smoothing resolutions, and performs robust outlier
  rejection across acquisition order. B) FlowMOP applies the same
  valid-parameter check, derives a candidate FSC-A threshold from each
  eligible parameter's positive events, and uses the median candidate as
  the final FSC-A gate. C) FlowMOP identifies doublets from inflection
  points in FSC-A/FSC-H and SSC-A/SSC-H ratio histograms, with a
  fixed-ratio fallback.''
  \end{quote}
\end{itemize}

\textbf{Location:} Figure 1 legend.

\subsection{Combined Comment 16 --- Abstract, Introduction, and Overall
Manuscript
Structure}\label{combined-comment-16--abstract-introduction-and-overall-manuscript-structure}

\textbf{Raised by:} Reviewer 2, general assessment and Comments 2 and 9.

\subsubsection{Reviewer 2 --- General Assessment (Manuscript
Structure)}\label{reviewer-2--general-assessment-manuscript-structure}

\begin{quote}
Also, the manuscript is not well structured. Careful edition by
experienced researcher is needed to increase the readability of the
manuscript.
\end{quote}

\subsubsection{Reviewer 2, Comment 2}\label{reviewer-2-comment-2}

\begin{quote}
The structure of abstract is kind of odd. The first paragraph should
focus on background instead of providing a conclusion.
\end{quote}

\subsubsection{Reviewer 2, Comment 9}\label{reviewer-2-comment-9}

\begin{quote}
The introduction is missing an end paragraph for summary.
\end{quote}

\subsubsection{Response}\label{response-15}

The manuscript has been considerably restructured and rewritten. The
Abstract now begins with the general preprocessing problem and
introduces FlowMOP directly. Temporal artifacts are then defined within
the methodological summary before the validation design and principal
findings are presented. We also added a concluding Introduction
paragraph that summarizes FlowMOP\textquotesingle s intended role and
previews the time-gating comparison, debris and doublet validation,
expert evaluation, and computational benchmark. Considering that the
paper has been considerably restructured, we hope that the new format
and writing address the reviewer\textquotesingle s concerns.

\subsubsection{Changes made}\label{changes-made-15}

\begin{itemize}
\item
  We reordered the Abstract so that it begins with the study context:

  \begin{quote}
  ``Flow cytometry now generates high-parameter datasets whose scale and
  variability challenge manual preprocessing, leading to subjectivity
  and poor reproducibility. Here, we introduce FlowMOP, a Python-native
  framework that automates three major preprocessing
  steps---time-gating, debris removal, and doublet exclusion.''
  \end{quote}
\item
  We added the following study summary to the end of the Introduction:

  \begin{quote}
  ``Here, we introduce FlowMOP, an automated preprocessing tool for
  time-gating, debris removal, and doublet exclusion. We compare
  time-gating performance with PeacoQC and FlowCut, evaluate debris and
  doublet removal using synthetic ground-truth datasets and expert
  comparison, and benchmark computational scalability.''
  \end{quote}
\item
  We reorganized and edited the manuscript so that the motivation,
  algorithmic design, technical benchmarks, biological validation,
  expert evaluation, and limitations are presented in a clearer
  progression.
\end{itemize}

\textbf{Location:} Abstract; final paragraph of the Introduction;
structure and writing throughout the manuscript.

\subsection{Associate Editor\textquotesingle s Closing
Request}\label{associate-editors-closing-request}

\begin{quote}
Finally, I ask that the authors also provide thorough and complete
responses to all comments raised by the other reviewers.
\end{quote}

\subsection{Comment Coverage Index}\label{comment-coverage-index}

\subsubsection{Associate Editor}\label{associate-editor}

{\def\LTcaptype{none} % do not increment counter
\begin{longtable}[]{@{}
  >{\raggedright\arraybackslash}p{(\linewidth - 2\tabcolsep) * \real{0.6400}}
  >{\raggedright\arraybackslash}p{(\linewidth - 2\tabcolsep) * \real{0.3200}}@{}}
\toprule\noalign{}
\begin{minipage}[b]{\linewidth}\raggedright
Associate Editor point
\end{minipage} & \begin{minipage}[b]{\linewidth}\raggedright
Response location
\end{minipage} \\
\midrule\noalign{}
\endhead
\bottomrule\noalign{}
\endlastfoot
General assessment & Associate Editor\textquotesingle s General
Assessment \\
Algorithmic motivation and theoretical justification --- existing
alternatives and design rationale & Combined Comment 1 \\
Algorithmic motivation and theoretical justification --- smoothing
mechanism & Combined Comment 4 \\
Algorithmic versus implementation advantages & Combined Comment 2 \\
Expert preference rankings, downstream biological impact, and
statistical interpretation & Combined Comments 3 and 6 \\
Smaller point --- dataset scale and complexity & Combined Comment 7 \\
Smaller point --- Bimix/Trimix and experimental relevance & Combined
Comment 9 \\
Request for complete responses & Associate Editor\textquotesingle s
Closing Request \\
\end{longtable}
}

\subsubsection{Reviewer 1}\label{reviewer-1}

{\def\LTcaptype{none} % do not increment counter
\begin{longtable}[]{@{}
  >{\raggedright\arraybackslash}p{(\linewidth - 2\tabcolsep) * \real{0.6400}}
  >{\raggedright\arraybackslash}p{(\linewidth - 2\tabcolsep) * \real{0.3200}}@{}}
\toprule\noalign{}
\begin{minipage}[b]{\linewidth}\raggedright
Reviewer 1 comment
\end{minipage} & \begin{minipage}[b]{\linewidth}\raggedright
Response location
\end{minipage} \\
\midrule\noalign{}
\endhead
\bottomrule\noalign{}
\endlastfoot
Comment 1 --- Technical infrastructure and performance claims & Combined
Comment 2 \\
Comment 2 --- Debris gating methodology & Combined Comment 10 \\
Comment 3 --- Refinement of human subjective rankings and dataset
breadth & Combined Comments 3 and 6 \\
Comment 4 --- Error costs and sensitivity trade-offs & Combined Comment
3 \\
Comment 5 --- Parametric sensitivity analysis & Combined Comment 4 \\
\end{longtable}
}

\subsubsection{Reviewer 2}\label{reviewer-2}

{\def\LTcaptype{none} % do not increment counter
\begin{longtable}[]{@{}
  >{\raggedright\arraybackslash}p{(\linewidth - 2\tabcolsep) * \real{0.6400}}
  >{\raggedright\arraybackslash}p{(\linewidth - 2\tabcolsep) * \real{0.3200}}@{}}
\toprule\noalign{}
\begin{minipage}[b]{\linewidth}\raggedright
Reviewer 2 comment
\end{minipage} & \begin{minipage}[b]{\linewidth}\raggedright
Response location
\end{minipage} \\
\midrule\noalign{}
\endhead
\bottomrule\noalign{}
\endlastfoot
General assessment & Combined Comments 3 and 16 \\
Comment 1 & Combined Comment 1 \\
Comment 2 & Combined Comment 16 \\
Comment 3 & Combined Comment 9 \\
Comment 4 & Combined Comment 9 \\
Comment 5 & Combined Comment 11 \\
Comment 6 & Combined Comment 1 \\
Comment 7 & Combined Comment 2 \\
Comment 8 & Combined Comment 1 \\
Comment 9 & Combined Comment 16 \\
Comment 10 & Combined Comment 3 \\
Comment 11 & Combined Comment 13 \\
Comment 12 & Combined Comment 13 \\
Comment 13 & Combined Comment 8 \\
Comment 14 & Combined Comment 9 \\
Comment 15 & Combined Comment 10 \\
Comment 16 & Combined Comment 6 \\
Comment 17 & Combined Comment 14 \\
Comment 18 & Combined Comment 5 \\
Comment 19 & Combined Comment 15 \\
Comment 20 & Combined Comment 4 \\
Comment 21 & Combined Comment 5 \\
Comment 22 & Combined Comment 10 \\
Comment 23 & Combined Comment 15 \\
Comment 24 & Combined Comment 10 \\
Comment 25 & Combined Comment 11 \\
Comment 26 & Combined Comment 11 \\
Comment 27 & Combined Comment 12 \\
Comment 28 & Combined Comment 8 \\
\end{longtable}
}

\subsection{Closing Statement}\label{closing-statement}

The revisions and explicit limitations provide a clearer account of the
manuscript\textquotesingle s evidentiary basis and practical
interpretation.

\end{document}
